% Appendix C: full construction + mechanism detail (moved from the main narrative).
\section{Benchmark construction and mechanism detail}
\label{app:detail}

\subsection{The attractor: post-SFT scores are benchmark-fixed, so ``stronger bases specialize more'' is arithmetic}
\label{sec:actI-coupling}

The cleanest way to see what Act~I does and does not show is to notice a regularity that runs underneath
\Cref{tab:primary}: \textbf{fine-tuning behaves like an attractor}. No matter where a checkpoint starts,
SFT pulls its ranking to nearly the \emph{same} score on a given benchmark. Represented-source AP after SFT
is $0.982 \pm 0.005$ across the four checkpoints; transfer AP after SFT is $0.807 \pm 0.024$ --- both far
tighter than the \emph{base} spread they came from ($0.178$ and $0.073$ respectively; the transfer variance
shrinks about \textbf{9-fold}). \Cref{fig:attractor} (left) shows the collapse: four widely spread base
scores, one narrow post-SFT band.

\begin{figure}[htbp]\centering
\includegraphics[width=0.99\linewidth]{fig_attractor.pdf}
\caption{\textbf{The fine-tuning attractor.} \emph{Left:} on the transfer benchmarks the four base
checkpoints span $0.79$--$0.94$ macro-AP, but after SFT they collapse into a narrow band around a
benchmark-fixed endpoint $A_T\approx0.81$. \emph{Right:} because every checkpoint lands near that same
endpoint, the paired change $\Delta=\text{SFT}-\text{base}$ is forced onto a line of slope $-1$ (green);
the four points sit on it (fitted slope $-1.05$, $R^2=0.91$). ``Stronger bases specialize more'' is
therefore arithmetic, not a discovered behavioral law.}
\label{fig:attractor}
\end{figure}

\paragraph{Why this reframes ``stronger bases specialize more.''} Once the endpoint is (approximately)
fixed, the change we plot is pure subtraction: $\Delta = \text{endpoint} - \text{base}$. A stronger base
then \emph{must} show a smaller $\Delta$ --- not because strong models are intrinsically fragile, but
because they start closer to the ceiling. This is testable. Regressing $\Delta$ on the
base score across the four checkpoints gives slope $-0.99$ (represented, $R^2=0.999$) and $-1.05$ (transfer,
$R^2=0.91$) --- the $-1$ the attractor \emph{predicts before any fitting}. One caveat on reading those
numbers: the represented $R^2=0.999$ is close to \emph{definitional}. When the endpoint band is as tight as
it is on represented sources ($\pm0.005$), $\Delta=\text{endpoint}-\text{base}$ is arithmetically almost a
slope-$-1$ line in the base score, so a near-perfect fit is nearly forced and is not independent
confirmation of anything. The transfer fit, where the endpoint band is several times wider, is the
load-bearing test --- and it still lands at slope $-1.05$. Either way the eye-catching ordering
(SmolLM2 $+0.040$ down to Qwen3-4B $-0.150$) carries no behavioral content beyond the endpoints themselves.

\paragraph{Where the real content lives.} Two places, both benchmark-owned rather than model-owned. First,
\emph{where the endpoints sit}: $A_T\approx0.81$ is largely a property of the \emph{(training manifest,
benchmark)} pair --- to the precision of the small residuals below, the base model has largely dropped out
of the equation, and the post-SFT score is \emph{largely determined by} the manifest and the benchmark
rather than by which checkpoint you started from. This is the mechanism behind the title's claim:
after SFT the score is largely a property of the (manifest, benchmark) pair, so a gain measured on that
pair says correspondingly little about any other. Second, the \emph{small residuals} around the
endpoint (at most $0.027$ on transfer): these measure how much of a base's identity survives the fine-tune,
and are arguably a more honest per-checkpoint quantity than $\Delta$ itself. LoRA only adds a low-rank
correction in the base's own features, so full base-independence is approximate by construction --- the
residual is exactly that approximation, made visible.

\begin{background}{Why the represented ceiling is $\approx0.98$, not $1.0$ --- and why that helps the thesis}
A ranker of the \emph{latent} truth generally cannot score AP $=1$ against slightly noisy \emph{observed}
labels (though it can rank the observed labels perfectly if the noise is systematic and learnable): under
random label noise a small rate $\eta$ of ambiguous or mislabeled rows lowers the expected achievable AP to
roughly $1-1.3\,\eta$ (the coefficient
$1.3$ is an order-of-magnitude rule of thumb, not derived here --- the exact factor depends on prevalence
and on where mislabels fall in the ranking). The observed common
ceiling $0.982$ is what you would see at $\eta\approx1.5\%$ label noise --- i.e.\ the ceiling would then be a
property of the benchmark's \emph{labels}, not of the models. This is a \emph{hypothesis}, not a measured
result: it is checkable by auditing the top-ranked ``false positives'' of the SFT guards --- under this
reading most should be mislabels --- but we have not run that audit, so the label-noise explanation of the
ceiling remains conjectural. Either way, the empirical fact we rely on is only that a common ceiling exists
and is shared across checkpoints; the benchmark, not the model, sets the top of the scale.
\end{background}

\begin{takeaway}{How much to trust the ``base strength'' ordering}
The apparent trend rests on \emph{four} checkpoints --- and because seeds within a checkpoint are highly
correlated (every Qwen3-4B seed is negative, every SmolLM2 seed but one positive), the effective sample is
$n=4$, not $20$. But the attractor turns the checkpoint-level ``specialize'' pattern into a \emph{prediction} from a
single constant: a checkpoint specializes exactly when its base transfer score exceeds the endpoint
($\text{base}_T>A_T\approx0.81$), which selects Qwen2.5, SmolLM3 and Qwen3-4B as specializers and leaves
SmolLM2 in uniform-gain --- the checkpoint-mean split we observe, and a $15/5$ seed count. The match is at
the checkpoint mean, not seed-perfect: two of the twenty seeds cross their checkpoint's boundary (Qwen2.5's
seed~42 gains $+0.008$; SmolLM2's seed~46 loses $-0.001$), so the constant predicts the aggregate split
rather than every individual seed. The sharper, falsifiable hypothesis to carry forward is
therefore \emph{endpoint invariance} itself: any new $1.5$--$4$B instruct checkpoint tuned with this recipe
should land near represented $0.98$ and transfer $0.81$ regardless of where its base started. One honest
limit on this extrapolation: our four checkpoints are only \emph{two} lineages (two Qwen, two SmolLM), so
what we have actually shown is endpoint invariance \emph{within} two families --- a prospective test
(roadmap item~7) should add unrelated lineages before the invariance is treated as recipe-general. What
Act~I establishes now is the qualitative claim, robust across all four checkpoints and $20$ seeds --- SFT
buys represented-source ranking and not, on average, transfer.
\end{takeaway}


\subsection{HMDA grounding: de-identified, banded fact sheets}
\label{sec:actIV-mortgage-hmda}

Realism is what makes the benchmark hard: a guard should face requests that read like genuine
mortgage-workflow traffic, not toy prompts. We ground each scenario in the public HMDA 2022 loan-level
snapshot \citep{hmda2022}, pulled from the FFIEC/CFPB Data Browser (the U.S.\ agencies that collect and
publish HMDA data). Each source record is reduced to a \emph{banded, de-identified fact sheet}: loan
purpose, occupancy, banded loan amount, banded income, banded LTV and DTI, action taken, denial reason,
and state --- and never an exact dollar amount, a census tract, or any identifier.

Three properties make this PII-safe by construction. \textbf{(i) Banding}: every exact figure is
replaced by a range (an income band rather than a dollar amount, likewise for LTV and DTI).
\textbf{(ii) Marginal sampling}: each field is drawn from its own marginal distribution over the bands
--- how often each band occurs \emph{on its own} --- rather than from the real joint combinations of
fields in any one record, so no actual borrower's row can be reassembled. \textbf{(iii) A build-time
assertion} that no emitted fact sheet reproduces a single source record verbatim; the frozen release
carries \code{contains\_real\_pii=false} with zero violations at freeze. Following
\citet{bowen2024mortgage}, the fair-lending and ability-to-repay cells are deliberately biased toward
\textbf{borderline, higher-risk} files (high DTI, high LTV, a prior denial) --- exactly where
underwriter discretion, and therefore bias, has room to operate.

\begin{background}{``Banded,'' ``marginal,'' and ``PII-safe'' --- in plain words}
\emph{Banded} means we never print an exact number; ``income \$92{,}450'' becomes ``income in the
\$75k--\$100k band.'' \emph{Marginal sampling} means we build a synthetic applicant field-by-field from
how common each band is by itself, not by copying a real person's combination of fields --- so even
though every band is realistic, the assembled applicant corresponds to no real borrower. \emph{PII-safe}
(personally identifiable information) then follows: with no exact figures, no tract, and no real
field-combination, there is nothing in a fact sheet that could re-identify anyone. The grounding buys
realism; the banding and marginal sampling buy privacy.
\end{background}

\subsection{The agentic construction pipeline}

\begin{figure}[htbp]\centering
\includegraphics[width=\linewidth]{pipeline.png}
\caption{\textbf{The agentic mortgage-benchmark construction pipeline.} Real HMDA records $\to$
de-identified banded fact sheets $\to$ the five-stage loop (plan / ground / generate /
adversarially mutate / rubric-bound judge), run once per coverage cell $\to$ accept-iff the judged
label matches the planned target $\to$ provenance and decontamination $\to$ the
\texttt{content\_family}-isolated split $\to$ the frozen \texttt{v1\_hmda2022} release, scored by the
evaluator. The two red nodes are the release's two stated defects, drawn rather than only written:
pairs were never a grouping key, so one protected pair crosses a split, and \texttt{private\_test} is
committed with text and already spent rather than sealed (\Cref{sec:actIV-mortgage-composition}). An
earlier version of this figure was a committed PNG with no source that still labelled that slice
``sealed''; it is now rendered from \code{figures/pipeline.dot} by the figure harness, so it cannot
drift from the prose again.}
\label{fig:pipeline}
\end{figure}
\label{sec:actIV-mortgage-pipeline}

Rows are authored and labeled by a five-stage agentic loop, summarized in \Cref{fig:pipeline}. The loop
is \emph{deterministic in structure} (the coverage cells it must fill are enumerated up front) but
\emph{stochastic in surface wording} (the LLM stages run at temperature $>0$), which is why the release
is a frozen artifact rather than a regenerable one.

\begin{enumerate}[leftmargin=1.6em,itemsep=2pt]
  \item \textbf{Planner.} Enumerates coverage cells as a product of quadrant $\times$ trap type
  $\times$ policy card $\times$ role $\times$ protected context, so every intended combination is
  targeted rather than left to chance.
  \item \textbf{HMDA-grounder.} Draws a de-identified banded fact sheet (\Cref{sec:actIV-mortgage-hmda})
  matching the planned cell.
  \item \textbf{Generator.} Authors the request in a mortgage-workflow \emph{role voice} --- applicant,
  loan officer, underwriter, processor, broker, or adversary --- so the phrasing matches who would
  plausibly send it.
  \item \textbf{Adversarial mutator.} Composes \emph{label-preserving} tactics that make a violation
  subtle and deniable without changing its ground truth: euphemism, coded proxy, business
  justification, and buried injection.
  \item \textbf{Rubric-bound judge.} Assigns $G$ and $D$ by reading the request against the policy
  cards (\Cref{sec:actIV-mortgage-cards}).
\end{enumerate}

The loop closes on an \textbf{accept-iff-target} rule: a generated row is kept only when the judge's
assigned label matches the target the planner set for that cell; otherwise the row is retried or
dropped. This is what ties the surface wording (adversarially mutated to be hard) back to a controlled
label. Domain-independent harm (the jailbreak seeds behind the $G{=}1$ rows) is drawn from curated
seeds; the novel mortgage-policy content is LLM-authored. Accepted rows then pass provenance and
decontamination, a \code{content\_family}-isolated train/dev/test split plus a committed 95-row extra
slice that is \emph{not} sealed and has already been spent (\Cref{sec:actIV}), and are published
as the frozen \code{v1\_hmda2022} release --- the object the reproducible harness scores.

Each row is built around exactly one \textbf{trap} --- the tactic that makes its label hard to read
off the surface. \code{direct} states the ask plainly; \code{business\_justified} wraps it in a
business rationale; \code{coded\_proxy} substitutes a stand-in (a neighborhood, a language) for a
protected group; \code{euphemism} softens the wording; \code{occupancy\_temptation} invites lying about
owner-occupancy; \code{buried\_injection} hides an instruction inside otherwise ordinary text;
\code{over\_refusal\_bait} looks alarming but is in fact benign, so a good guard should \emph{not} flag
it; \code{benign\_info} is a plain safe request; and \code{minimal\_pair} is the protected-class
counterfactual of \Cref{sec:actIV-mortgage-pairs}.

\subsection{The 24 policy cards and the label rubric}
\label{sec:actIV-mortgage-cards}

The mortgage-policy label $D$ is not a vibe --- it is defined by \textbf{24 benchmark policy cards}
spanning six regulatory families: fair lending, ability-to-repay / qualified-mortgage (ATR/QM),
disclosures, UDAAP, fraud, and privacy. Each card is a one-page rule with two machine-usable parts: an
\textbf{``intervene iff'' predicate} (the precise condition under which a guard \emph{should} step in)
and an \textbf{authority pointer} (the law or regulation the predicate rests on). The judge reads a
request against the applicable cards and emits $G\in\{\text{safe},\text{unsafe}\}$ and
$D\in\{\text{allow},\text{intervene}\}$; the composed $\mathrm{final}$ of \Cref{eq:mortgage-final}
follows mechanically, together with an \emph{action lattice} (the ordered menu of responses, from
allow through warn to block) and a \emph{severity} grade.

The honesty caveat is load-bearing and stated plainly: \textbf{these labels are policy-card-consistent,
assigned by an LLM judge --- they are not SME-adjudicated.} The judge is run rubric-bound against the
applicable cards at temperature~$0$ with a three-sample majority vote; all 24 cards remain unsigned, and
judge agreement is measured as \emph{self-consistency} (re-running the same judge and checking it agrees
with itself), not as a human inter-rater study. Two transparency gaps we flag rather than resolve here: the
frozen artifact does not record the exact judge \emph{model and version}, so its
possible model-family overlap with an evaluated checkpoint cannot be ruled out as a source of correlated
blind spots; and per-label self-consistency \emph{counts} were not exported, so self-consistency is asserted,
not tabulated.

\begin{background}{Policy cards, ``SME,'' and Fleiss-\texorpdfstring{$\kappa$}{kappa} --- in plain words}
A \emph{policy card} turns a slice of mortgage law into something a program can check: a plain
``intervene when \ldots'' condition plus a citation to the rule it comes from. An \emph{SME} is a
subject-matter expert --- here, a compliance lawyer --- and \emph{Fleiss-$\kappa$} is a standard number
between $0$ and $1$ for how much several \emph{independent human} raters agree beyond chance, the usual
gold standard for a labeled benchmark. We do not yet have that. What we report instead is
\emph{self-consistency}: the same LLM judge, re-run, agreeing with itself --- a much weaker guarantee.
That gap is precisely why every result in this section is a diagnostic, not a certified fair-lending
finding.
\end{background}

